\documentclass[11pt]{amsart}
\usepackage[T1]{fontenc}
\usepackage{lmodern}
\usepackage{amsmath,amssymb,amsthm,mathtools}
\usepackage{microtype}
\usepackage{xcolor}
\usepackage[colorlinks=true,linkcolor=blue!45!black,citecolor=blue!45!black,urlcolor=blue!45!black]{hyperref}

\newtheorem{theorem}{Theorem}[section]
\newtheorem{lemma}[theorem]{Lemma}
\newtheorem{proposition}[theorem]{Proposition}
\newtheorem{corollary}[theorem]{Corollary}
\theoremstyle{definition}
\newtheorem{definition}[theorem]{Definition}
\newtheorem{remark}[theorem]{Remark}

\newcommand{\ii}{\mathrm{i}}
\newcommand{\C}{\mathbb{C}}
\newcommand{\R}{\mathbb{R}}
\newcommand{\DeltaA}{\mathit{\Delta}}

\title[Atiyah--Sutcliffe linear independence for two pyramid families]
{Atiyah--Sutcliffe linear independence for centered rectangular and rhombic pyramids}
\author{Zijian Zeng}
\address{Institute of Computer Science and Digital Innovation, UCSI University,
Kuala Lumpur 56000, Malaysia}
\email{zijianzeng@foxmail.com}
\email{1002266693@ucsiuniversity.edu.my}
\date{August 30, 2026}
\subjclass[2020]{51M15, 15A03}
\keywords{Atiyah--Sutcliffe conjecture, configurations of points, binary forms,
symbolic factorization}
\hypersetup{
  pdftitle={Atiyah--Sutcliffe linear independence for centered rectangular and rhombic pyramids},
  pdfauthor={Zijian Zeng},
  pdfsubject={Atiyah--Sutcliffe Conjecture 1 for two five-point families}
}

\begin{document}

\begin{abstract}
We prove Atiyah--Sutcliffe Conjecture~1 for two nonplanar five-point
families. In each family four points form either an arbitrary rectangle or
an arbitrary rhombus, and the fifth point lies off their plane on the normal
line through the center. After rationalizing the apex directions, the two
binary-quartic determinants admit explicit factorizations into strictly
positive quantities, up to a nonzero constant phase. The identities are
verified by two exact computations: coefficient convolution and direct
evaluation at five projective points followed by a Vandermonde determinant.
We also record a boundary lemma reducing a nondegenerate one-point escape to
the known four-point theorem. These results do not prove the conjecture for
arbitrary five points in three-space.
\end{abstract}

\maketitle

\section{Introduction}

For distinct points $x_1,\ldots,x_n\in\R^3$, Atiyah and Sutcliffe associate
to every $x_i$ a binary form of degree $n-1$, using the directions from
$x_i$ to the other points. Their Conjecture~1 asserts that the resulting
$n$ forms are linearly independent, or equivalently that the associated
determinant is nonzero \cite{AtiyahSutcliffe2002}.

The conjecture is known for arbitrary configurations with $n\leq4$;
the four-point case is due to Eastwood and Norbury
\cite{EastwoodNorbury2001}, with a later proof by Malkoun
\cite{Malkoun2023}. A number of special configurations are also known.
In particular, Djokovi\'c proved the conjecture when some points lie on an
axis and the remaining points form a regular polygon in a perpendicular
plane with center on that axis \cite{Djokovic2002}. For five points this
contains a centered square pyramid. Further results and determinant
formulas appear in \cite{MazurPetrenko2012,Malkoun2019}; a recent account
still separates the unrestricted $n\leq4$ result from the five-point
frontier \cite{Malkoun2026}.

We treat two centered but nonregular deformations of the square pyramid.

\begin{theorem}\label{thm:intro}
Atiyah--Sutcliffe Conjecture~1 holds for each of the following five-point
configurations:
\begin{enumerate}
\item four vertices of an arbitrary nondegenerate rectangle and a fifth
point off its plane on the normal line through its center;
\item four vertices of an arbitrary nondegenerate rhombus and a fifth point
off its plane on the normal line through its center.
\end{enumerate}
\end{theorem}

Each part is a two-real-parameter family up to Euclidean similarities. The
square occurs in both parts and is already covered by Djokovi\'c's theorem;
the content here is that no equality of the two base dimensions is needed.
The proof is constructive: Sections~\ref{sec:rectangle} and
\ref{sec:rhombus} give the exact determinant factorizations. Section
\ref{sec:verification} describes two independent exact ways to check them.

\section{Binary forms and a convenient lift}\label{sec:binary}

We fix conventions because individual lifts affect the displayed
determinant, although never its zero set. The Hopf map is
\[
h(a,b)=\bigl(2\operatorname{Re}(a\overline b),
2\operatorname{Im}(a\overline b),|a|^2-|b|^2\bigr).
\]
For an ordered pair $i\ne j$, choose a lift $(a_{ij},b_{ij})$ of the unit
direction $(x_j-x_i)/\lVert x_j-x_i\rVert$ and set
\[
L_{ij}(U,V)=a_{ij}V-b_{ij}U,
\qquad p_i(U,V)=\prod_{j\ne i}L_{ij}(U,V).
\]
Writing
\[
p_i(U,V)=\sum_{k=0}^4c_{ki}U^{4-k}V^k,
\qquad \DeltaA=\det(c_{ki})_{0\leq k\leq4,\ 1\leq i\leq5},
\]
Conjecture~1 is the assertion $\DeltaA\ne0$. Multiplying any $L_{ij}$ by a
nonzero scalar only rescales one column, so nonvanishing is independent of
these choices.

We use the following unnormalized lift throughout.

\begin{lemma}\label{lem:lift}
Let $(X,Y,Z)\ne0$ have length $r$ and suppose $(X,Y,Z)$ is not the north-pole
vector $(0,0,r)$. Then
\[
\ell_{(X,Y,Z)}(U,V)=-(r-Z)U+(X+\ii Y)V                     \tag{2.1}
\]
is a valid Atiyah linear factor for the direction $(X,Y,Z)/r$.
\end{lemma}

\begin{proof}
Take $(a,b)=(X+\ii Y,r-Z)$. Its squared norm is $2r(r-Z)$, and a direct
calculation gives
\[
h(a,b)=2(r-Z)(X,Y,Z).
\]
After normalization this is $(X,Y,Z)/r$, while $aV-bU$ is (2.1).
\end{proof}

No displacement used below is vertical: the apex is above the center rather
than above a base vertex. Thus Lemma~\ref{lem:lift} applies to every factor.

\section{Centered rectangular pyramids}\label{sec:rectangle}

Up to similarity and reflection, every configuration in the first part of
Theorem~\ref{thm:intro} has the form
\begin{equation}\label{eq:rectangle-points}
\begin{aligned}
x_1&=(1,A,0),&x_2&=(-1,A,0),\\
x_3&=(-1,-A,0),&x_4&=(1,-A,0),\qquad x_5=(0,0,H),
\end{aligned}
\end{equation}
where $A,H>0$. Define
\[
S=\sqrt{1+A^2},\qquad R=\sqrt{S^2+H^2},\qquad
w=\frac{R-H}{S}=\frac{S}{R+H}.
\]
Then $S>0$, $0<w<1$, and
\begin{equation}\label{eq:w-inverse}
H=\frac{S(1-w^2)}{2w},\qquad R=\frac{S(1+w^2)}{2w}.
\end{equation}

For a base vertex write $z_i=(x_i)_1+\ii(x_i)_2$. Formula (2.1) gives
\begin{align*}
L_{i5}&=-wS\,U-z_iV &&(1\leq i\leq4),\\
L_{5j}&=-\frac{S}{w}U+z_jV &&(1\leq j\leq4),
\end{align*}
while a base-to-base factor is
$-r_{ij}U+(z_j-z_i)V$, with $r_{ij}\in\{2,2A,2S\}$.

Set
\[
E=w^6+w^5+w+1
\]
and
\begin{equation}\label{eq:rectangle-Q}
\begin{split}
Q={}&(A^4+1)(w^3+1)^2
 +SE(A^3+A^2+A+1)\\
 &+2(w^5+w)(A^3+A)+2A^2(w^3-1)^2.
\end{split}
\end{equation}

\begin{proposition}\label{prop:rectangle-factor}
For the determinant convention of Section~\ref{sec:binary},
\begin{equation}\label{eq:rectangle-factor}
\begin{split}
w^4\DeltaA={}&-2^{16}A^4S^5(A^2+AS+S+1)\\
&\quad\cdot\bigl(S(1+w^2)+2w(A+1)\bigr)Q.
\end{split}
\end{equation}
The identity holds in
$\mathbb Z[\ii,A,S,w]/(S^2-A^2-1)$.
\end{proposition}

\begin{proof}
Expand each product of four displayed linear factors in the ordered basis
$U^4,U^3V,\ldots,V^4$. Multiplying the fifth column by $w^4$ clears its
denominators. Exact determinant expansion and reduction by
$S^2-A^2-1$ gives the right side of \eqref{eq:rectangle-factor}.
Section~\ref{sec:verification} gives two reproducible checks that do not share
the same determinant representation.
\end{proof}

\begin{proof}[Proof of Theorem~\ref{thm:intro}(1)]
Every displayed factor on the right of \eqref{eq:rectangle-factor}, apart
from the leading minus sign, is strictly positive for $A,S,w>0$. In
particular, all four groups in \eqref{eq:rectangle-Q} are nonnegative and the
first three are strictly positive. Since $w\ne0$, we obtain
$\DeltaA\ne0$.
\end{proof}

\section{Centered rhombic pyramids}\label{sec:rhombus}

The diagonals of a nondegenerate rhombus are perpendicular. Hence, up to
similarity and reflection, every configuration in the second part of
Theorem~\ref{thm:intro} can be written
\begin{equation}\label{eq:rhombus-points}
\begin{gathered}
x_1=(1,0,0),\quad x_2=(0,A,0),\quad x_3=(-1,0,0),\\
x_4=(0,-A,0),\qquad x_5=(0,0,H),\qquad A,H>0.
\end{gathered}
\end{equation}
Put
\[
S=\sqrt{1+A^2},\qquad R=\sqrt{1+H^2},\qquad
T=\sqrt{A^2+H^2}
\]
and introduce two positive apex parameters
\begin{equation}\label{eq:uv}
u=R-H=\frac1{R+H},\qquad
v=\frac{T-H}{A}=\frac{A}{T+H}.
\end{equation}
In fact $0<u,v<1$. With $z_i=(x_i)_1+\ii(x_i)_2$, put
$(d_i,t_i)=(1,u)$ for $i=1,3$ and $(d_i,t_i)=(A,v)$ for
$i=2,4$. The apex factors are
\[
L_{i5}=-d_it_iU-z_iV,\qquad
L_{5i}=-\frac{d_i}{t_i}U+z_iV.
\]
The six base-pair lengths are $S,S,S,S,2,2A$.

Define
\begin{equation}\label{eq:rhombus-P}
\begin{split}
P={}&A^3uv+A^3v+A^2Suv+A^2Su+A^2Sv+A^2S+A^2uv+A^2u\\
 &+AS^2u+AS^2+2ASu+2ASv+Auv+Av\\
 &+S^2v+S^2+Suv+Su+Sv+S+uv+u
\end{split}
\end{equation}
and
\begin{equation}\label{eq:rhombus-Q}
\begin{split}
Q={}&A^4(u+v)+2A^3S(v+1)\\
 &+A^2S^2\bigl(2u^3v^3+(u+v)(u-v)^2+2\bigr)\\
 &+2A^2S(u+1)+2A^2(u+v)\\
 &+2AS^3(u^3v^3+u^3v^2)+2AS(v+1)\\
 &+S^4(u^3v^2+u^2v^3)+2S^3(u^3v^3+u^2v^3)\\
 &+2S^2u^3v^3+S^2\bigl((u+v)(u-v)^2+2\bigr)\\
 &+2S(u+1)+(u+v).
\end{split}
\end{equation}

\begin{proposition}\label{prop:rhombus-factor}
For the determinant convention of Section~\ref{sec:binary},
\begin{equation}\label{eq:rhombus-factor}
u^2v^2\DeltaA=2^7\ii A^6SPQ.
\end{equation}
This is an identity in the polynomial ring
$\mathbb Z[\ii,A,S,u,v]$, before imposing any geometric relation among the
four variables.
\end{proposition}

\begin{proof}
Expand the quartics from the factors above. Multiplication of the fifth
column by $u^2v^2$ clears every denominator. Exact determinant expansion
gives \eqref{eq:rhombus-factor}. In obtaining the displayed positive form of
$Q$, the only signed four-term group is rewritten using
\begin{equation}\label{eq:square-regroup}
u^3+v^3-u^2v-uv^2=(u+v)(u-v)^2.
\end{equation}
Independent exact checks are described in Section~\ref{sec:verification}.
\end{proof}

\begin{proof}[Proof of Theorem~\ref{thm:intro}(2)]
Every monomial in \eqref{eq:rhombus-P} is strictly positive. Every summand
in \eqref{eq:rhombus-Q} is nonnegative, and several are strictly positive.
Thus $P,Q>0$. Since $A,S,u,v>0$, equation
\eqref{eq:rhombus-factor} implies $\DeltaA\ne0$.
\end{proof}

\section{Exact verification}\label{sec:verification}

The factorizations were verified in SymPy~1.14.0 by two separate exact
determinant constructions.

\subsection{Coefficient construction}
For each linear factor $c_0U+c_1V$, the first verifier updates the coefficient
list by convolution. It then forms the $5\times5$ coefficient matrix and
uses an exact domain-matrix determinant. The unreduced rectangular
determinant has $231$ terms over $\mathbb Z[\ii,A,S,w]$; reduction modulo
$S^2-A^2-1$ gives \eqref{eq:rectangle-factor}. The cleared rhombic
determinant has $363$ terms over $\mathbb Z[\ii,A,S,u,v]$ and equals
\eqref{eq:rhombus-factor} identically.

\subsection{Projective-value construction}
The second verifier does not form quartic coefficients. For
$t\in\{0,1,-1,\ii,-\ii\}$, it directly multiplies the four values
$c_0+c_1t$ for every column. If $C$ is the coefficient matrix and $M$ is
the resulting value matrix, then
\[
M=VC,\qquad V=(t_j^k)_{0\leq j,k\leq4},\qquad \det V=-16\ii.
\]
Consequently $\det C=\det M/(-16\ii)$. This computation independently
reproduces both claimed identities with exact Gaussian-integer arithmetic.

The verification scripts, located in the archive's \texttt{code} directory,
are
\begin{itemize}
\item \texttt{certify\_rectangular\_pyramid.py},
\item \texttt{verify\_rectangular\_pyramid\_identity.py},
\item \texttt{certify\_rhombic\_pyramid.py}, and
\item \texttt{verify\_rhombic\_pyramid\_identity.py}.
\end{itemize}
Each program reconstructs its determinant from the displayed linear factors
and exits unsuccessfully if the difference from the claimed expression is
nonzero. The short positivity arguments following the identities are not
delegated to numerical computation.

\section{A one-point escape boundary}\label{sec:escape}

The following observation is useful for attempts at the unrestricted
five-point problem.

\begin{lemma}\label{lem:escape}
Fix four distinct bounded points $x_1,\ldots,x_4\in\R^3$. If a fifth point
escapes to infinity in a fixed direction $e\in S^2$, then the limiting five
binary quartics in the directional compactification are linearly independent.
Therefore, sufficiently far along that ray, the five-point determinant is
nonzero.
\end{lemma}

\begin{proof}
Let $q_1,\ldots,q_4$ be the four Atiyah cubics of the bounded configuration.
All directions from a bounded point to the escaping point tend to $e$, and
all reverse directions tend to $-e$. Up to nonzero rescaling, the limiting
quartics are
\[
L_eq_1,\quad L_eq_2,\quad L_eq_3,\quad L_eq_4,\quad L_{-e}^4.
\]
The antipodal factors $L_e,L_{-e}$ are independent. A binary change of
variables sends them to $U,V$. In the basis
$U^4,U^3V,U^2V^2,UV^3,V^4$, the five-point coefficient matrix is then block
triangular: its upper left block is the four-point coefficient matrix of the
$q_i$, and its last diagonal entry is nonzero. The upper block is nonsingular
by the four-point theorem \cite{EastwoodNorbury2001}. The limiting
determinant is therefore nonzero, and continuity gives the final assertion.
\end{proof}

The lemma does not address collisions, degeneration of the four-point
cluster, or simultaneous escape of several clusters.

\section{Limitations}

Theorem~\ref{thm:intro} covers two two-parameter similarity families inside
the eight-dimensional similarity moduli of generic ordered five-point
configurations. It does not establish Conjecture~1 for arbitrary five
points. In particular, the rectangle and rhombus factorizations cannot be
interpolated without proof to an arbitrary parallelogram base.

As a diagnostic only, a deterministic floating-point search tested $70{,}000$
random, near-planar, near-collinear, clustered, and small-integer
configurations, together with three global-optimization runs. It produced no
candidate zero. This negative observation is not used anywhere in the
proofs. The exact verification described in Section~\ref{sec:verification}
uses no floating-point arithmetic.

\section*{Data and code availability}

The accompanying source archive contains the LaTeX manuscript, the exact
verification scripts, machine-readable certificate summaries, and pinned
Python package versions. The core verification commands are documented in
the included README file.

\newpage
\bibliographystyle{amsplain}
\bibliography{references}

\end{document}
